\documentclass[10pt,twocolumn]{article}

%% Packages
\usepackage[margin=0.9in]{geometry}
\usepackage{amsmath,amssymb,amsthm}
\usepackage{graphicx}
\usepackage{booktabs}
\usepackage{hyperref}
\usepackage{xcolor}
\usepackage{microtype}
\usepackage{cite}
\usepackage{enumitem}
\usepackage{float}
\usepackage{caption}
\usepackage{subcaption}
\usepackage{multirow}

%% Theorem environments
\newtheorem{theorem}{Theorem}
\newtheorem{lemma}[theorem]{Lemma}
\newtheorem{corollary}[theorem]{Corollary}
\newtheorem{definition}{Definition}
\newtheorem{proposition}[theorem]{Proposition}
\theoremstyle{remark}
\newtheorem{remark}{Remark}

%% Macros
\newcommand{\vrfi}{\textsc{VRFI}}
\newcommand{\vlr}{\textsc{VLR}}
\newcommand{\mct}{\textsc{MCT}}

\title{\textbf{Learning to Schedule: A Value-Aware Policy Discovery\\
Experiment with Provable Queue-Global Necessity}}

\author{Dev Upadhyay\\
M.S.\ Computer Science, University of California Santa Cruz\\
Graduate Capstone Project, 2024--2025\\
\texttt{dupadhyay@ucsc.edu}}

\date{March 2026}

\begin{document}

\maketitle

%% ============================================================
\begin{abstract}
%% ============================================================

CPU scheduling's efficiency-fairness tension has resisted formal
characterization. We train a two-head attention DQN on Alibaba 2018
cluster traces with a value-delta reward, where each task carries
urgency curve $V_i(d) = \mathrm{base}_i \cdot
\max(\mathrm{floor}_i,\, e^{-d/\tau_i})$, aggregating
$V_i(d+q) - V_i(d)$ across all runnable processes.

We prove the \emph{Quantum Scoring Monotonicity Lemma}: any
per-task local scorer based on value-drop objectives is strictly
monotone in quantum size, making non-degenerate quantum selection
impossible without queue-global information---confirmed by three
ablations. W10C, our two-head attention DQN, achieves
$21.71 \pm 0.64$\,s mean completion time (replicated across 3
independent seeds, CV\,=\,2.9\%), beating Round Robin with zero
starvation and lowest slowdown variance. The best-observed single
checkpoint reaches 17.23\,s (N\,=\,1500 evaluation episodes,
original checkpoint \texttt{dqn\_w10c.npz}).

We introduce the \emph{Value-Rate Fairness Index (\vrfi)}, proving
a case exists where JFI\,=\,1.0 and SDV\,=\,0 yet
\vrfi\,=\,0.641---exposing value-rate inequality invisible to prior
metrics. A 2$\times$2 ablation isolates reward formulation and burst
observability as independently load-bearing: value-delta reward
outperforms composite reward by 9.1\,s MCT when burst is available,
while removing burst costs only 4.4\,s---establishing that the oracle
information gap is modest on realistic workloads.

A preference-conditioned variant (W14-$\omega$) beats MLFQ on both
throughput (20.60\,s vs.\ 21.59\,s MCT) and starvation prevention
(32.5\% vs.\ 36.0\%) at $\omega_s=0.7$, using no oracle information
and requiring no retraining to shift operating modes.

\end{abstract}

%% ============================================================
\section{Introduction}
%% ============================================================

Modern CPU schedulers operate on what is immediately observable:
processes in the ready queue, their recent CPU usage, and priority
class. Linux's Completely Fair Scheduler (CFS) distributes CPU time
proportionally to virtual runtime. But fairness over what is visible
is not the same as minimizing the time tasks spend waiting to become
visible. A short-burst process that arrives while the CPU is occupied
by a long-burst process experiences a waiting time that CFS cannot
reduce by design---the problem is not in its policy but in its horizon.

We ask: can an agent learn a scheduling policy \emph{without behavioral
cloning or reference policy demonstrations}---no supervision, no target
policy to imitate, no handcrafted priority rules---purely from the
signal of how long processes have been alive? This is a policy
discovery problem, not a CFS improvement problem. We do not propose
to replace CFS. We ask whether RL can, in a controlled environment,
rediscover approximations to the known optimal policy, and whether
failure modes reveal architectural requirements for generalizing
across arbitrary workloads.

\subsection*{Contributions}

\begin{enumerate}[leftmargin=*, label=\arabic*., itemsep=2pt]

\item \textbf{Systematic failure mode documentation} across four
architectures: tabular discretization collapse, DQN position-indexed
gradient contamination, spurious load-adaptive quantum artifact, and
attention entropy collapse. Each is diagnosed to a specific
architectural cause and resolved or left explicitly open.

\item \textbf{W10C two-head attention DQN}: $21.71\pm0.64$\,s MCT
replicated across 3 seeds (CV\,=\,2.9\%), beating Round Robin by
$\sim$2\,s. Best observed checkpoint: 17.23\,s (N\,=\,1500 eval).
Spontaneous head specialization---context/survey and starvation
monitor---emerges without supervision.

\item \textbf{Quantum Scoring Monotonicity Lemma and three ablations}:
formal proof that value-drop local scorers $f_{\mathrm{rate}}$ and
$f_{\mathrm{total}}$ always select $q_{\min}$ and $q_{\max}$
respectively, for all $\tau>0$, $d\geq0$. Three ablations confirm
this empirically on Alibaba-derived trace data.

\item \textbf{Gradient alignment and $\tau$ calibration}: identifies a
3--4 order-of-magnitude disconnect between default $\tau$ values
(25\,s/100\,s) and trace-scale delays (medians 5,279\,s/38,665\,s),
producing near-ties in 93--100\% of ordering decisions regardless of
calibration---a structural consequence of trace composition.

\item \textbf{\vrfi\ and JFI/SDV falsification}: a new metric
$1 - \mathrm{CV}(\vlr)$ capturing value-curve fairness invisible to
Jain's Index and slowdown variance. Key falsification: identical
delays yield JFI\,=\,1.0, SDV\,=\,0, yet \vrfi\,=\,0.641.

\item \textbf{2$\times$2 ablation} over reward structure and burst
observability: reward formulation is more load-bearing than state
observability. Removing value-delta reward (W11d) degrades MCT to
26.35\,s; removing oracle burst (W12) degrades MCT to 21.67\,s.
Oracle information cost is 4.4\,s on the 2018 trace.

\item \textbf{W14-$\omega$ preference-conditioned agent}: beats MLFQ
simultaneously on MCT (20.60\,s vs.\ 21.59\,s) and starvation
(32.5\% vs.\ 36.0\%) at $\omega_s=0.7$. Non-monotonic Pareto
frontier establishes that optimal fairness requires retaining
throughput pressure.

\item \textbf{W15 variable-N training}: retraining on Poisson arrival
streams recovers the MCT advantage (mean $-1.40$\,s vs.\ MLFQ across
3 seeds), confirming the attention mechanism generalizes to variable
queue depths when trained on them.

\end{enumerate}

%% ============================================================
\section{Background}
%% ============================================================

\subsection{CFS and the Scheduling Horizon Problem}

CFS uses a red-black tree ordered by virtual runtime to select the
next runnable process, targeting equal CPU share among active tasks.
It makes no attempt to prefer shorter jobs---the key property of SRPT,
the preemptive optimal policy for minimizing mean turnaround time when
burst lengths are known. Since Linux kernel 6.6, EEVDF has replaced
CFS as the default Linux scheduler~\cite{eevdf}. EEVDF assigns each
task a virtual deadline based on weight and requested scheduling slice,
preventing tasks from receiving CPU before they are entitled to it.
The value-aware scheduler proposed here differs from both: scheduling
decisions are driven by explicit per-task value curves rather than
virtual runtime accounting.

\subsection{RL for Scheduling}

KernelOracle~\cite{kernelOracle} applies supervised learning to
predict CFS's next selection---it does not learn a new policy.
ALPS~\cite{alps} simulates SRPT on historical workloads and enforces
priorities via eBPF; it achieves 57.2\% latency reduction on
serverless traces. Both approaches are anchored to a reference policy.
RLScheduler~\cite{rlscheduler} learns HPC batch job scheduling from
rich per-job state including historical runtime statistics.
Double DQN scheduling~\cite{doubleDQN} achieves strong performance on
mixed CPU-I/O workloads. SmartOS~\cite{smartOS} applies RL to joint
multi-resource allocation. The common thread: domain-informed state
design with scheduler engineering knowledge. This work deliberately
restricts state to features requiring no domain knowledge to isolate
what reward signal alone can teach.

\subsection{Queue-Global Information in Production Systems}

A 2026 public preview of cloud batch scheduler fairness introduced a
mechanism requiring server-side backlog tracking that \emph{cannot}
be implemented using per-task metadata alone---the backlog count is a
property of the queue, not any individual task. Section~\ref{sec:theory}
formalizes this: any value-drop scoring function operating on per-task
features is provably unable to make optimal quantum selections, because
both natural scoring objectives are monotone in quantum size. Breaking
this degeneracy requires queue-global information---precisely the
server-side backlog such systems track.

%% ============================================================
\section{Theoretical Framework}
\label{sec:theory}
%% ============================================================

\subsection{Value Curve Formulation}

Each task $i$ carries an exponential value curve:
\begin{equation}
  V_i(d) = \mathrm{base}_i \cdot
  \max\!\bigl(\mathrm{floor}_i,\; e^{-d/\tau_i}\bigr)
  \label{eq:vcurve}
\end{equation}
where $d \geq 0$ is accumulated queueing delay, $\tau_i > 0$ is the
urgency decay constant. Two canonical profiles emerge from the Alibaba
trace: \textbf{steep} ($\tau \approx 25$\,s, floor$\,=0.2$, modeling
interactive/IO tasks) and \textbf{smooth} ($\tau \approx 100$\,s,
floor$\,=0$, modeling batch tasks). Note: curve type assignment
(steep/smooth) is a human modeling choice encoding task urgency; it
is not inferred from the trace.

\subsection{Potential Surface}

Define $\Phi(S) = \sum_i V_i(d_i)$ over queue state $S$. The gradient:
\begin{equation}
  \nabla_i \Phi = \frac{\mathrm{base}_i}{\tau_i}\,e^{-d_i/\tau_i}
  \quad\text{(non-plateau)}; \quad = 0 \quad\text{(plateau)}
\end{equation}

A locally optimal scheduler selects the task with largest
$|\nabla_i\Phi|$. This is inherently queue-global: selecting task $i$
depends not on $i$'s features alone, but on whether any co-queued task
$j$ has $|\nabla_j\Phi| > |\nabla_i\Phi|$.

\subsection{Quantum Scoring Monotonicity}
\label{sec:lemma}

Define two natural quantum scoring functions for task $i$ given
quantum $q$:
\begin{align}
  f_{\mathrm{rate}}(q) &= \frac{V_i(d_i) - V_i(d_i+q)}{q}
  \label{eq:frate}\\
  f_{\mathrm{total}}(q) &= V_i(d_i) - V_i(d_i+q)
  \label{eq:ftotal}
\end{align}

\begin{lemma}[Quantum Scoring Monotonicity]
\label{lem:mono}
For any $\tau > 0$, $\mathrm{base} > 0$, $\mathrm{floor} \in [0,1)$,
$d \geq 0$:
\begin{enumerate}[(a)]
  \item $f_{\mathrm{rate}}(q)$ is strictly decreasing in $q > 0$.
  \item $f_{\mathrm{total}}(q)$ is strictly increasing in $q > 0$
    such that $d+q$ is below the plateau.
\end{enumerate}
\end{lemma}

\begin{proof}
\textbf{(a):} Let $h(q) = (1-e^{-q/\tau})/q$. Then
$f_{\mathrm{rate}}(q) = \mathrm{base}\cdot e^{-d/\tau}\cdot h(q)$.
Differentiating: $h'(q) = [(q/\tau+1)e^{-q/\tau}-1]/q^2$.
Setting $u=q/\tau>0$, the numerator $(u+1)e^{-u}-1<0$ for all $u>0$
since $e^u>1+u$ by strict convexity. Therefore $h'(q)<0$. \hfill$\square$

\textbf{(b):} Below the plateau, $f_{\mathrm{total}}(q) =
\mathrm{base}\cdot e^{-d/\tau}\cdot(1-e^{-q/\tau})$, which is
strictly increasing since $1-e^{-q/\tau}$ is strictly increasing
in $q$ for $\tau>0$. \hfill$\square$

\textbf{Plateau case:} When $d_i$ exceeds the plateau,
$f_{\mathrm{total}}(q)=0$ identically---degenerate because the task
carries no further scheduling signal.
\end{proof}

\begin{corollary}[Queue-Global Necessity for Value-Drop Scorers]
\label{cor:global}
No scoring function based solely on value-drop---operating on
$(\mathrm{base}_i, \mathrm{floor}_i, \tau_i, d_i, q)$ via either
$f_{\mathrm{rate}}$ or $f_{\mathrm{total}}$---can produce mixed
quantum selection across $\{q_{\min}, q_{\mathrm{med}}, q_{\max}\}$.
$f_{\mathrm{rate}}$ always selects $q_{\min}$; $f_{\mathrm{total}}$
always selects $q_{\max}$. This corollary applies specifically to
these natural value-preservation objectives; other local rules
(e.g., threshold rules on $d/\tau$) could in principle produce mixed
selection while remaining local. Queue-global information is necessary
to escape the degenerate poles of $f_{\mathrm{rate}}$ and
$f_{\mathrm{total}}$.
\end{corollary}

\subsection{Queue-Global Resolution}

The impossibility results motivate maintaining queue-global state at
each decision step. The two-head attention mechanism (Section
\ref{sec:arch}) is one instantiation: one head specializes in
competitor context (queue state) and the other in per-task scoring.
Section \ref{sec:theory-confirms} confirms this empirically.

%% ============================================================
\section{Method}
\label{sec:method}
%% ============================================================

\subsection{Environment}

\textbf{Synthetic (Weeks 1--9).} Five processes per episode. Arrivals
sampled from $\{0,2,5,8,10\}$\,ms; bursts from $U[1,60]$\,ms.
Actions: 15 (5 processes $\times$ 3 quantum tiers: 1\,ms, 5\,ms,
20\,ms). Invalid actions masked at every step.

\textbf{Real-trace (Weeks 8--15).} Alibaba 2018 cluster trace,
\texttt{batch\_task} table (14.3M tasks). Filtering: $0 <
\mathrm{duration} \leq 47.0$\,s (p75 cutoff). Split: 80/20
train/test (8.18M / 2.04M tasks). Quantum tiers: 0.5\,s, 2.0\,s,
8.0\,s for real-trace agents (15 actions: 5 processes $\times$ 3
quanta).

\subsection{Reward}

\textbf{MCT reward (W1--W10C).}
$R(t) = -n_{\mathrm{active}}(t) \times q_{\mathrm{actual}}(t)$,
where $n_{\mathrm{active}}$ is the count of arrived, incomplete
processes and $q_{\mathrm{actual}} = \min(q, \mathrm{burst\_remaining})$.
By Little's Law, minimizing $\sum_t n_{\mathrm{active}} \times
q_{\mathrm{actual}}$ is equivalent to minimizing total turnaround
time. No completion bonuses or shaped incentives.

\textbf{Value-delta reward (real-trace variants).}
\begin{equation}
  R = \frac{1}{20}\sum_{i \in \mathrm{runnable}}
  \bigl[V_i(d_i+q) - V_i(d_i)\bigr]
  \label{eq:vdreward}
\end{equation}
Dense (fires every step), implicitly encodes urgency and completion
pressure.

\textbf{W14-$\omega$ composite reward.}
$R = \omega_{\mathrm{mct}} \cdot R_{\mathrm{vd}} +
\omega_s \cdot R_{\mathrm{ss}}$, where
$R_{\mathrm{ss}} = -\sum_i \max(0, t_{\mathrm{starve},i} - 50)/50\,/N$
is a proportional starvation signal.

\subsection{State Representation}

\textbf{Oracle disclosure.} W10C includes true remaining burst time
as feature \texttt{off+0} (\texttt{remaining\_norm})---an oracle
assumption unavailable in production. W12 and W15 replace this with
observable time-based features: time\_in\_queue, wait\_time,
time\_since\_last\_exec, urgency\_norm, cpu\_norm, mem\_norm,
arrived\_flag (35-dim, 5 processes $\times$ 7 features). The cost
of this substitution is quantified in the 2$\times$2 ablation
(Sections \ref{sec:2x2a}--\ref{sec:2x2b}): $\sim$4.4\,s MCT on the
2018 trace.

\subsection{Architecture}
\label{sec:arch}

\begin{figure}[t]
  \centering
  \includegraphics[width=\columnwidth]{fig3_architecture.pdf}
  \caption{W14-$\omega$ architecture: 2-head attention DQN with
  FiLM $\omega$-conditioning at query modulation (pre-attention:
  $Q_{\mathrm{cond}} = Q \times (1+\omega_s)$) and post-attention
  context ($c_{\mathrm{cond}} = c\times(1+\omega_s)+\omega_{\mathrm{mct}}$).
  MLP input: $[\mathbf{e}_i(7)\,\|\,\mathbf{c}_i(16)\,\|\,\omega_s(1)]
  = 24$-dim.}
  \label{fig:arch}
\end{figure}

\textbf{Attention DQN (W10C).} 2 heads, $D_{\mathrm{head}}=8$.
Per-head attention over the 5-process encoding, with scaled
dot-product attention. Output projected to 16-dim context vector.
Per-candidate scoring MLP: $[e_i(7)\,\|\,q_t(1)] \to 64 \to 32 \to 1$,
where $q_t$ is the normalized quantum tier. Total: 4,145 parameters.

\textbf{Quantum selection note.} In W10C, the quantum tier index
enters the MLP as an explicit feature, giving genuinely independent
Q-values for each (task, quantum) pair: the MLP can assign different
scores to (task 2, $q=0.5$\,s) and (task 2, $q=8.0$\,s). In W15
(PyTorch), one task-level score is computed per candidate and tiled
across all 3 quantum tiers---quantum selection is uniform among valid
actions for the chosen task. W10C's quantum selection is learned;
W15's is not.

\textbf{FiLM conditioning (W14-$\omega$).} Preference scalar
$\omega_s \in [0,1]$ modulates computation at two points:
\begin{align}
  Q_{\mathrm{cond}} &= Q \cdot (1+\omega_s)
  \quad\text{(pre-attention query modulation)}\\
  c_{\mathrm{cond}} &= c\cdot(1+\omega_s) + \omega_{\mathrm{mct}}
  \quad\text{(post-attention context)}
\end{align}
Fixed scaling (not learned)---$\omega$ acts as a control signal.

\textbf{Omega sampling (W14-$\omega$).}
Episodes 0--8000: $\omega_s \sim U(0,1)$.
Episodes 8000+: $\omega_s \sim \mathrm{Beta}(0.5,0.5)$ (approximated
via $\omega_s = \sin^2(u\pi/2)$, $u \sim U[0,1]$). Replay buffer
stores $\omega_s$ per transition; never resampled at training time.

\subsection{Baselines}

\textbf{SRPT oracle:} optimal mean completion time, requires true
burst times. \textbf{Round Robin:} fixed-quantum preemptive, no
ordering. \textbf{MLFQ:} 3 queues, quanta $\{4,8,16\}$\,s, aging
threshold 50\,s. \textbf{CFS-lite:} vruntime with 1\,s granularity,
10\,s target latency. \textbf{EEVDF-lite:} CFS-lite extended with
per-task eligible times.

%% ============================================================
\section{Results}
\label{sec:results}
%% ============================================================

\subsection{W10C: Primary Result}

\begin{figure}[t]
  \centering
  \includegraphics[width=\columnwidth]{fig2_learning_curve.pdf}
  \caption{W10C training progression. Round Robin (21.31\,s) and
  SRPT oracle (15.20\,s) shown as references. Gray band:
  $21.71\pm0.64$\,s across 3 fresh training seeds (CV\,=\,3.0\%).}
  \label{fig:learning}
\end{figure}

Table~\ref{tab:w10c} reports the full agent progression.
\textbf{W10C (replicated mean): $21.71\pm0.64$\,s} across 3
independent seeds (CV\,=\,3.0\%), beating Round Robin by
$\sim$2\,s. Best observed checkpoint: $17.23\pm11.38$\,s
(N\,=\,1500 episodes, \texttt{dqn\_w10c.npz}$^\dagger$). See
Figure~\ref{fig:learning}.

\begin{table}[h]
\centering
\caption{Agent progression on Alibaba 2018 trace. All values
reported on filtered test split. $^\dagger$Original checkpoint,
N\,=\,1500 eval; replicated fresh training: $21.71\pm0.64$\,s
(3 seeds, CV\,=\,3.0\%).}
\label{tab:w10c}
\small
\begin{tabular}{lrr}
\toprule
Agent & MCT (s) & vs RR \\
\midrule
SRPT oracle & 15.20 & --- \\
W10C orig.\ checkpoint$^\dagger$ & 17.23$\pm$11.38 & $-$4.08 \\
\textbf{W10C replicated (3 seeds)} & \textbf{21.71$\pm$0.64} & $\sim$$-$2.0 \\
W9 1-head & 19.22$\pm$11.58 & $-$2.09 \\
Round Robin & 21.31$\pm$14.03 & baseline \\
\bottomrule
\end{tabular}
\end{table}

Two attention heads specialize spontaneously without supervision:
Head 1 attends broadly across the queue (context/survey); Head 2
consistently monitors the longest-waiting competitor (starvation
prevention). SRPT agreement: 71.9\% overall, 69.9\% on steep tasks,
30.3\% on smooth tasks.

\subsection{Quantum Degeneracy Ablation}
\label{sec:theory-confirms}

Three ablations confirm Lemma~\ref{lem:mono} (Table~\ref{tab:abl}):

\begin{table}[h]
\centering
\caption{Quantum degeneracy ablation.}
\label{tab:abl}
\small
\begin{tabular}{lrrr}
\toprule
Variant & $q_{\min}$\% & $q_{\max}$\% & Total Value \\
\midrule
\texttt{quantum\_only} ($f_{\mathrm{rate}}$) & 100 & 0 & 250.55 \\
\texttt{quantum\_only\_fixed} ($f_{\mathrm{total}}$) & 0 & 100 & 159.21 \\
\texttt{ordering\_only} ($q_{\mathrm{med}}$) & --- & --- & 15.37 \\
\textbf{W10C (attention)} & mixed & mixed & 241.92 \\
\bottomrule
\end{tabular}
\end{table}

Both monotone scorers collapse to degenerate poles. The attention DQN
produces mixed quantum selection---evidence that queue-global context
breaks the degeneracy.

\subsection{Gradient Alignment and $\tau$ Calibration}

Default $\tau$ values (25\,s, 100\,s) are miscalibrated: median
delays are 5,279\,s and 38,665\,s. At these delays,
$e^{-5279/25}\approx 10^{-92}$---numerically zero. The value gradient
is zero; 93--100\% of ordering decisions become near-ties regardless
of policy. Calibrated values ($\tau_{\text{steep}}=900$\,s,
$\tau_{\text{smooth}}=800$\,s) restore gradient signal but near-ties
persist due to trace composition (92\% smooth tasks).

\subsection{Value-Rate Fairness Index (\vrfi)}
\label{sec:vrfi}

\begin{figure}[t]
  \centering
  \includegraphics[width=\columnwidth]{fig4_vrfi_falsification.pdf}
  \caption{\vrfi\ falsification (Example D): 10 tasks with identical
  delays ($d=50$\,s) yield JFI\,=\,1.0, SDV\,=\,0 under existing
  metrics, yet \vrfi\,=\,0.641 correctly detects that steep tasks
  lose value at twice the rate of smooth tasks.}
  \label{fig:vrfi}
\end{figure}

\begin{definition}[\vlr\ and \vrfi]
For task $i$ with delay $d_i$:
\begin{equation}
  \vlr_i = \frac{V_i(0) - V_i(d_i)}{d_i},
  \quad \lim_{d_i\to 0}\vlr_i = \frac{\mathrm{base}_i}{\tau_i}
  \label{eq:vlr}
\end{equation}
\begin{equation}
  \vrfi = 1 - \frac{\sigma(\vlr)}{\mu(\vlr)} = 1 - \mathrm{CV}(\vlr)
  \label{eq:vrfi}
\end{equation}
$\vrfi\in(-\infty,1]$: value 1 = equal value-loss rates; $\to-\infty$
= concentrated loss.
\end{definition}

\textbf{Falsification.} Ten tasks, identical delay $d=50$\,s: 5 steep
($\tau=25$\,s) and 5 smooth ($\tau=100$\,s). Then JFI\,=\,1.0
(equal turnarounds), SDV\,=\,0 (equal slowdowns), yet steep tasks
have $\vlr\approx0.016$ vs.\ smooth $\vlr\approx0.008$---a $2\times$
disparity. \vrfi\,=\,0.641 correctly identifies the inequality.
See Figure~\ref{fig:vrfi}.

Table~\ref{tab:fair} reports five-policy evaluation. Value-aware
policy achieves rank sum 17 vs.\ MLFQ's 28, zero starvation, lowest
global SDV.

\begin{table}[h]
\centering
\caption{Five-policy fairness evaluation (300-task slice).}
\label{tab:fair}
\small
\begin{tabular}{lrrrr}
\toprule
Policy & JFI & SDV & Starv & \vrfi \\
\midrule
value\_aware & 0.934 & 566 & \textbf{0} & 0.712 \\
mlfq & 0.943 & 2,193 & 0 & 0.420 \\
cfs\_lite & 0.916 & 6,138 & 1 & --- \\
rr & \textbf{0.953} & 15,819 & 5 & --- \\
fcfs & 0.774 & 12.7M & 27 & --- \\
\bottomrule
\end{tabular}
\end{table}

\subsection{2$\times$2 Ablation: Reward $\times$ Observability}
\label{sec:2x2a}
\label{sec:2x2b}

\begin{table}[h]
\centering
\caption{2$\times$2 ablation (Alibaba 2018, mean$\pm$std, 3 seeds).
$^*$W10C original checkpoint N\,=\,1500; fresh training:
$21.71\pm0.64$\,s.}
\label{tab:2x2}
\small
\begin{tabular}{llrr}
\toprule
Condition & Agent & MCT (s) & SRPT\,(\%) \\
\midrule
Neither removed & W10C & $17.23\pm11.38^*$ & 45.7$\pm$3.0 \\
Burst removed & W12 & $21.67\pm0.42$ & 56.6 \\
Reward changed & W11d & $26.35\pm1.93$ & 40.6 \\
Both changed & W11/W11b & $21.05\pm0.72$ & 54.7 \\
\bottomrule
\end{tabular}
\end{table}

Three findings: (1)~\textbf{Reward dominates when burst is available}:
W10C vs.\ W11d (same state, different reward) shows 9.1\,s gap.
(2)~\textbf{Burst costs 4.44\,s}: W10C vs.\ W12 (same reward, no
burst). (3)~\textbf{Without burst, reward barely matters}: W12
vs.\ W11/W11b differ by 0.6\,s---state representation is the
binding constraint.

Furthermore, W11d (26.35\,s) is \emph{worse} than W11/W11b
(21.05\,s) despite having more state information. Burst time actively
interferes with composite reward learning.

\textbf{Core insight.} The value-delta reward implicitly unifies
efficiency and urgency-weighted fairness: a task early in its decay
generates a large negative delta, approximating SRPT-like ordering
as an emergent consequence. Decomposing this into additive components
breaks the coupling and causes policy collapse (Section \ref{sec:pbrs}).

\subsection{PBRS Lambda Tuning}
\label{sec:pbrs}

PBRS with $\varphi(s) = -\lambda\cdot\max_i(t_{\mathrm{starve},i})$
tested at $\lambda\in\{0.01,0.05,0.10,0.25,0.50\}$: starvation
floor 46--57\% at all $\lambda$, unmoved. No $\lambda$ achieves
MCT\,$<21.59$\,s and starvation\,$<36.0\%$ simultaneously. MLFQ's
aging rule is a deterministic constraint; PBRS produces probabilistic
tendencies that the value-delta signal overrides at every tested
$\lambda$.

\subsection{Preference-Conditioned Agent (W14-$\omega$)}

\begin{figure}[t]
  \centering
  \includegraphics[width=\columnwidth]{fig1_pareto_frontier.pdf}
  \caption{W14-$\omega$ Pareto frontier as $\omega_s$ varies 0$\to$1.
  Shaded region dominates MLFQ on both metrics. $\omega_s=0.7$
  (gold star) recommended: MCT\,=\,20.60\,s, starvation\,=\,32.5\%.
  Non-monotonic frontier: pure starvation optimization
  ($\omega_s=1.0$, 38.5\%) is worse than balanced ($\omega_s=0.7$).}
  \label{fig:pareto}
\end{figure}

Gradient norm ratio 0.92--1.12 throughout 30k episodes---no
policy collapse. Table~\ref{tab:omega} shows the full frontier.

\begin{table}[h]
\centering
\caption{W14-$\omega$ Pareto frontier vs.\ MLFQ.}
\label{tab:omega}
\small
\begin{tabular}{rrrr}
\toprule
$\omega_s$ & MCT (s) & Starv\,(\%) & vs MLFQ \\
\midrule
0.0 & 20.51 & 44.0 & MCT\,$\checkmark$, Starv\,$\times$ \\
0.5 & 20.73 & 38.0 & MCT\,$\checkmark$, Starv\,$\times$ \\
\textbf{0.6} & \textbf{20.49} & \textbf{35.0} & \textbf{Both}\,$\checkmark$ \\
\textbf{0.7} & \textbf{20.60} & \textbf{32.5} & \textbf{Both}\,$\checkmark$ \\
\textbf{0.8} & \textbf{20.43} & \textbf{34.5} & \textbf{Both}\,$\checkmark$ \\
0.9 & 20.45 & 38.5 & MCT\,$\checkmark$, Starv\,$\times$ \\
1.0 & 20.36 & 38.5 & MCT\,$\checkmark$, Starv\,$\times$ \\
\midrule
MLFQ & 21.59 & 36.0 & baseline \\
\bottomrule
\end{tabular}
\end{table}

$\omega_s\in\{0.6,0.7,0.8\}$ simultaneously dominate MLFQ.
Recommended: $\omega_s=0.7$ (lowest starvation at 32.5\%).
The \textbf{non-monotonic frontier}---starvation minimized at
$\omega_s=0.7$, not $\omega_s=1.0$---establishes that optimal
fairness requires retaining throughput pressure.

\subsection{Variable-N Analysis and W15}

W14-$\omega$ trained on fixed $N=5$ loses to MLFQ at $\rho\geq0.7$
(OOD\%\,$>13.8\%$). Root cause: 5-slot truncation removes queue
visibility when $N>5$; MLFQ's aging rule operates over the full
queue regardless of size.

W15 retrains W14-$\omega$'s architecture on Poisson arrivals
($\lambda\sim U(0.01,0.08)$, 300 completions/episode), presenting
the 5 longest-waiting tasks in the attention window when $N>5$.
Training required four stability fixes: TD target clipping $[-50,50]$,
Huber loss, frequent target sync (every 5 episodes), LR\,=\,$10^{-4}$.

\begin{table}[h]
\centering
\caption{W15 results: variable-N training vs.\ MLFQ.}
\label{tab:w15}
\small
\begin{tabular}{lrr}
\toprule
Agent & MCT (mean$\pm$std, s) & Starv\,(\%) \\
\midrule
MLFQ baseline & $22.06\pm13.36$ & 37.5 \\
W15 seed 42 & $21.97\pm13.81$ & 40.0 \\
W15 seed 123 & $20.57\pm11.96$ & 47.5 \\
W15 seed 456 & $\mathbf{19.43\pm11.83}$ & 42.0 \\
W15 mean$\pm$std & $20.66\pm1.03$ & 43.2 \\
\bottomrule
\end{tabular}
\end{table}

W15 beats MLFQ on MCT across all 3 seeds (best: $-2.63$\,s; mean:
$-1.40$\,s). Starvation (40--48\%) remains above MLFQ (38\%)---
consistent with Section~\ref{sec:pbrs}: reward shaping alone cannot
enforce hard starvation guarantees.

%% ============================================================
\section{Discussion}
%% ============================================================

\subsection{What Value-Delta Reward Encodes}

The value-delta reward appears to be a fairness signal. In practice,
it implicitly encodes SRPT-like ordering: a task early in decay
generates a large negative delta, pushing the agent to complete it
quickly. The W11 ablation series confirmed this by decomposing the
signal---when split into additive components, both objectives degrade.
The signal's power comes from its unity.

\subsection{The Starvation Problem}

Starvation (40--53\% under various agents vs.\ MLFQ's 36--38\%) is
structural, not a tuning failure. MLFQ's aging rule is a deterministic
constraint. Even W14-$\omega$'s best setting ($\omega_s=0.7$, 32.5\%)
achieves only a 3.5\,pp advantage over MLFQ---a margin that warrants
caution in interpretation.

\subsection{Limitations}

\begin{enumerate}[itemsep=2pt]
\item \textbf{Environment scale.} Five processes, 300-task episodes,
  single core. Production schedulers handle millions of concurrent
  processes.
\item \textbf{Burst time.} W10C uses true remaining burst time
  (oracle). W12/W15 remove this at $\sim$4.44\,s MCT cost.
\item \textbf{Value curve assignment.} Steep/smooth curve assignment
  is a human modeling choice. In multi-tenant systems, $\tau$
  assignment could become implicit prioritization.
\item \textbf{Statistical rigor.} Three seeds is a minimum; 5--10
  seeds would provide stronger statistical guarantees for the W14-$\omega$
  and W15 results.
\end{enumerate}

\subsection{Ethical Considerations}

Value-aware scheduling raises equity questions beyond the technical.
Curve assignment (steep/smooth) follows task class---a heuristic that
may disadvantage batch workloads belonging to resource-constrained
users. The $\tau$ parameter directly controls urgency; in multi-tenant
systems it could enable implicit prioritization. The choice of fairness
definition (\vrfi\ vs.\ JFI) is a value judgment that should involve
stakeholder input in production deployments.

%% ============================================================
\section{Conclusion}
%% ============================================================

We trained a two-head attention DQN on the Alibaba 2018 cluster trace,
progressing through fifteen agent versions.

The \textbf{Quantum Scoring Monotonicity Lemma} proves formally that
$f_{\mathrm{rate}}$ and $f_{\mathrm{total}}$ are degenerate local
scorers---queue-global information is necessary to escape their poles.
This result connects to production systems that independently reached
the same architectural conclusion.

\textbf{\vrfi} fills a measurable gap: identical delays can yield
JFI\,=\,1.0, SDV\,=\,0, yet \vrfi\,=\,0.641---a concrete case where
value-rate inequality is invisible to prior metrics.

The \textbf{2$\times$2 ablation} establishes that reward structure and
burst observability are independently load-bearing. The oracle
information gap is 4.44\,s on the 2018 trace---modest relative to
earlier experiments on synthetic distributions.

\textbf{W14-$\omega$} beats MLFQ on both MCT and starvation at
$\omega_s=0.7$, with a single model adaptable at runtime via
preference vector $\omega$. The non-monotonic Pareto frontier
establishes that optimal fairness requires throughput pressure.

\textbf{W15} confirms MCT advantage generalizes to variable-N Poisson
arrival streams (mean $-1.40$\,s vs.\ MLFQ across 3 seeds).

\subsection*{Future Work}

(1) Preference-conditioned variable-N training combining W14-$\omega$
and W15. (2) Learned burst prediction to close the 4.44\,s oracle
gap. (3) Full transformer architecture scaling to variable $N$.
(4) LLM-based policy discovery using open-source models.

%% ============================================================
\bibliographystyle{plain}
\begin{thebibliography}{9}

\bibitem{kernelOracle}
  S.~Y.~Kahu.
  \newblock KernelOracle: Predicting the Linux Scheduler's Next Move
  with Deep Learning.
  \newblock {\em arXiv:2505.15213}, 2025.

\bibitem{alps}
  Y.~Fu, R.~Shi, H.~Wang, S.~Chen, and Y.~Cheng.
  \newblock ALPS: An Adaptive Learning, Priority OS Scheduler for
  Serverless Functions.
  \newblock {\em USENIX ATC}, 2024.

\bibitem{rlscheduler}
  D.~Zhang, D.~Dai, Y.~He, F.~S.~Bao, and B.~Xie.
  \newblock RLScheduler: An Automated HPC Batch Job Scheduler Using
  Reinforcement Learning.
  \newblock {\em SC}, 2020. doi:10.1109/SC41405.2020.00035.

\bibitem{doubleDQN}
  X.~Sun et~al.
  \newblock Dynamic OS Scheduling Using Double DQN.
  \newblock {\em arXiv:2503.23659}, 2025.

\bibitem{smartOS}
  S.~Goodarzy et~al.
  \newblock SmartOS: Towards Automated Learning and User-Adaptive
  Resource Allocation in OS.
  \newblock {\em APSys}, 2021. doi:10.1145/3476886.3477519.

\bibitem{littlesLaw}
  J.~D.~C.~Little.
  \newblock A Proof for the Queuing Formula: $L = \lambda W$.
  \newblock {\em Operations Research}, 9(3):383--387, 1961.

\bibitem{jainFairness}
  R.~Jain, D.~Chiu, and W.~Hawe.
  \newblock A Quantitative Measure of Fairness and Discrimination for
  Resource Allocation in Shared Computer Systems.
  \newblock DEC TR-301, 1984.

\bibitem{alibaba}
  Alibaba Group.
  \newblock Alibaba Cluster Trace Program, 2018.
  \newblock \url{https://github.com/alibaba/clusterdata}.

\bibitem{eevdf}
  T.~Natvig.
  \newblock EEVDF Scheduler.
  \newblock {\em Linux Kernel Mailing List}, 2023.
  \newblock \url{https://lwn.net/Articles/925371/}.

\end{thebibliography}

\end{document}
